% decay_betaproc: the three beta processes seen inside the nucleus
% (Vol III Ch 2, fig:decay-betaproc)
% beta-minus: n -> p + e- + anti-nu ;  beta-plus: p -> n + e+ + nu ;
% electron capture: p + e-(K shell) -> n + nu, then an outer electron fills the
% K vacancy and a characteristic X-ray is emitted.
% Charge and lepton number L_e balance in each: e-(+1) with anti-nu(-1);
% e+(-1) with nu(+1); captured e-(+1) becomes nu(+1).  A is unchanged.
% Colours: shared `proton' and `neutron' discs; electron cPrim, positron cCrim,
% neutrinos cTeal dashed, photon: shared wavy `photon' style (bookstyle legend).
\begin{tikzpicture}[
  nucleon/.style={circle,minimum size=6.5mm,inner sep=0pt},
  P/.style={nucleon,proton},
  N/.style={nucleon,neutron},
  nucl/.style={thin line,draw=cSteel,fill=cMoon!45},
  lepton/.style={-smallhead,construction},
  nutrack/.style={-smallhead,construction,cTeal,dash pattern=on 3pt off 1.5pt},
  xray/.style={-smallhead,construction,photon}]
% ----------------------------------------------------------- beta minus
\node[note] at (2.0,2.6) {$\beta^-$ decay};
\draw[nucl] (0.7,1.0) circle[radius=0.75]; \node[N] at (0.7,1.0) {$n$};
\node[small label,text=cSteel] at (0.7,0.0) {nucleus};
\draw[vector,cInk] (1.6,1.0) -- (2.3,1.0);
\draw[nucl] (3.1,1.0) circle[radius=0.55]; \node[P] at (3.1,1.0) {$p$};
\draw[lepton,cPrim] (3.55,1.3) -- (4.15,1.95) node[above,cPrim] {$e^-$};
\draw[nutrack] (3.55,0.7) -- (4.15,0.05) node[below,cTeal] {$\bar\nu_e$};
\node[small label,align=center] at (2.0,-0.85) {$n\to p+e^-+\bar\nu_e$\\$Z\to Z+1$, $A$ unchanged};
% ------------------------------------------------------------ beta plus
\begin{scope}[xshift=5.3cm]
\node[note] at (2.0,2.6) {$\beta^+$ decay};
\draw[nucl] (0.7,1.0) circle[radius=0.75]; \node[P] at (0.7,1.0) {$p$};
\node[small label,text=cSteel] at (0.7,0.0) {nucleus};
\draw[vector,cInk] (1.6,1.0) -- (2.3,1.0);
\draw[nucl] (3.1,1.0) circle[radius=0.55]; \node[N] at (3.1,1.0) {$n$};
\draw[lepton,cCrim] (3.55,1.3) -- (4.15,1.95) node[above,cCrim] {$e^+$};
\draw[nutrack] (3.55,0.7) -- (4.15,0.05) node[below,cTeal] {$\nu_e$};
\node[small label,align=center] at (2.0,-0.85) {$p\to n+e^++\nu_e$\\$Z\to Z-1$, $A$ unchanged};
\end{scope}
% ------------------------------------------------------ electron capture
\begin{scope}[xshift=10.7cm]
\node[note] at (1.3,2.6) {electron capture};
\draw[thin line,cTeal,dash pattern=on 1pt off 1.5pt] (1.2,1.0) circle[radius=0.95];
\node[small label,text=cTeal] at (1.2,2.15) {K shell};
\draw[nucl] (1.2,1.0) circle[radius=0.42]; \node[P] at (1.2,1.0) {$p$};
\node[dot,cPrim,minimum size=4pt] at (2.15,1.0) {};
\node[cPrim] at (2.45,1.3) {$e^-$};
\draw[-smallhead,thin line,cPrim] (2.1,0.9) to[bend left=30] (1.68,0.92);
\draw[nutrack] (0.85,0.68) -- (0.15,0.02) node[below,cTeal] {$\nu_e$};
\draw[xray] (2.05,0.55) -- (2.75,-0.05) node[below,small label,text=cAcc] {X-ray};
\node[small label,align=center] at (1.3,-0.85) {$p+e^-\to n+\nu_e$\\$Z\to Z-1$, $A$ unchanged};
\end{scope}
\end{tikzpicture}
